\documentclass[11pt,a4paper]{article}
\usepackage[margin=27mm]{geometry}
\usepackage{amsmath,amssymb,amsthm,booktabs,fontspec,unicode-math,hyperref}
\setmainfont{STIXGeneral}\setmathfont{STIX Two Math}
\setlength{\emergencystretch}{2em}
\newtheorem{theorem}{Theorem}\newtheorem{proposition}[theorem]{Proposition}
\newcommand{\PSL}{\operatorname{PSL}}\newcommand{\PGL}{\operatorname{PGL}}
\title{The chiral projective classification from three common interfaces}
\author{The human\_compression project}\date{11 October 2026}
\begin{document}\maketitle
\begin{abstract}
We recover the full principal classification and constructions of
Leemans--Moerenhout--O'Reilly-Regueiro, in the exact arXiv-v1 edition.
The common theory supplies anchored rank descent, two affine ends and
one exceptional trace quadratic. These give all the source's prime-field
branches and pointed class counts, with the omitted residue restored.
The project's earlier accepted primitive-square and translated-subfield
proofs also reproduce the subsequent completion for odd composite degrees.
The whole shared paper, complete finite premises and retained classical
inputs are explicit dependencies; no partial branch is presented as the
whole source.
\end{abstract}

Write $L(q)=\PSL_2(q)$ and $P(q)=\PGL_2(q)$, $q=p^d\ge4$.
For an odd prime $p\ne5$ with $5$ a square, put
\[
 D_4=\{1+\sqrt5,1-\sqrt5\},\quad
 D_{55}=\{(7+5\sqrt5)/2,(7-5\sqrt5)/2\},\quad
 D_{35}=\{3+2\sqrt5,3-2\sqrt5\}.
\]
For such a list $D$, let $s(D)$ count its nonzero square entries in
$F_p$. If $5$ is not a square set all three counts to zero; also set
$s(D_4)=0$ unless $2$ is a square.

\begin{theorem}\label{app:classification}
Neither $L(q)$ nor $P(q)$ is the full group of a chiral polytope of
rank three or at least five. In rank four:
\begin{enumerate}
\item $P(q)$ occurs exactly for $q>4$.
\item For even $q$, the same holds for $L(q)=P(q)$.
For $q\equiv1\pmod4$, $L(q)$ occurs exactly for $q>9$.
\item If $q=p^d\equiv3\pmod4$ with $d>1$ a prime power, $L(q)$ does not occur.
\item If $q=p\equiv3\pmod4$ is prime, every type, up to duality,
is $(5,3,4)$, $(5,3,5)$ or $(3,5,3)$, and their numbers of ordered
isomorphism classes are respectively $2s(D_4)$, $2s(D_{55})$, $2s(D_{35})$.
In particular existence is equivalent to at least one count being positive.
\item If $q=p^d\equiv3\pmod4$ and $d>1$ is not a prime power, $L(q)$ occurs.
This final clause reproduces the project's previously accepted completion,
not a result asserted as proved in the source edition.
\end{enumerate}
\end{theorem}
\begin{proof}
The simultaneous-reversal theorem of [U] excludes rank three; its anchored
involution theorem and subfield envelope exclude all ranks at least five.
The pointed construction and the complete $A_5,A_6$ exclusions give the
first two rank-four clauses. The nonsingular trace extension formula
puts a marking over a quadratic extension of the compositum of the two
proper end fields; when $d>1$ is an odd prime power this compositum is
proper and the quadratic extension is absent. This is the common
odd-degree exclusion.

For prime $p\equiv3\pmod4$, neither end fixes a projective point.
The common subgroup reduction and full exceptional local-pair census
leave exactly the three displayed types. The trace-extension and
collision theorems then give all the stated counts, generation,
intersection conditions and chirality. In particular $p=11$ contributes
no $(3,5,3)$ class, while $p=19$ contributes two $(5,3,5)$ classes.
These small conclusions are exact finite premises, not assumptions
borrowed from the source's table. Finally separate the prime-power
factors of a non-prime-power odd $d$ into coprime $e_1,e_2>1$ and apply
the common coprime-subfield construction. Every part has both its
necessity and sufficiency supplied in the paid common paper.
\end{proof}

For direct comparison with source Theorem~1(4), the five existence
conditions in its prime branch are
\[
\begin{array}{c|l}
(a)&p\equiv11,19\pmod{20},\quad
 p\bmod11\in\{1,3,4,5,9\},\quad D_{35}\text{ both squares};\\
(b)&p\equiv11,19\pmod{20},\quad
 p\bmod11\in\{2,6,7,8,10\};\\
(c)&p\equiv11,19\pmod{20},\quad
 p\bmod19\in\{1,4,5,6,7,9,11,16,17\},\quad D_{55}\text{ both squares};\\
(d)&p\equiv11,19\pmod{20},\quad
 p\bmod19\in\{2,3,8,10,12,13,14,15,18\};\\
(e)&p\equiv31,39\pmod{40}.
\end{array}
\]
The restored residue $10$ in (b) is present in the source's Theorem~11
but missing from its introductory Theorem~1. The common discriminant
products $-11,-19,-4$ and quadratic reciprocity prove precisely these
conditions. The case $p=19$ is already included by (b). As in the
source, both enantiomorphic ordered classes are counted when distinct;
the displayed $(5,3,4)$ counts also hold for its dual $(4,3,5)$.
These are the source proof's ordered generating-tuple classes under
$\operatorname{Aut}(G)$; no further identification by changing the
base flag to the opposite flag orbit is imposed, and no separate
unbased-poset census is inferred from them.

\section{All pointed ends, not only one existence example}
The following includes the full counts of source Theorems~4 and~5.
For a positive integer $k>2$ require
\[
 k\mid(q-1),\qquad k\nmid p^e-1,\quad k\nmid p^e+1
       \quad(0<e<d).
\]
For odd $q$ also require $k$ even. If the group is $L(q)$ require
$k\mid(q-1)/2$; if it is $P(q)\ne L(q)$ require
$k\nmid(q-1)/2$. In characteristic two these two conditions are
replaced by equality $P(q)=L(q)$, with no even-$k$ condition.
There are exactly $\varphi(k)/d$ ordered chiral markings whose two
ends each fix a projective point. Their common type is
$(k,k,k)$ in characteristic two or when $4\mid k$, and
$(k/2,k,k/2)$ when $k\equiv2\pmod4$.

Here is the complete source marking, with an element $t$ of order $k$:
\[
 a=\begin{pmatrix}-t^{-1}&1\\0&1\end{pmatrix},\quad
 b=\begin{pmatrix}t&0\\0&1\end{pmatrix},\quad
 c=\begin{pmatrix}1&0\\1+t&-t\end{pmatrix}.
\]
The common pointed theorem proves both ends are
$F_q^+\rtimes\langle-1,t\rangle$, their intersection is the common
diagonal group, and they generate the whole specified projective group.
It proves the normal form is exhaustive and the only identifications
are Frobenius on $t$. Thus the construction gives all classes and every
claimed count. The condition excluding a Frobenius inversion of $t$
excludes a regular marking, not necessarily a proper generated subgroup.
For example $q=25$, $|t|=6$ gives all of $L(25)$, with end orders $150$
and intersection six, but its Frobenius reversal makes it regular.
This corrects that proof route in the source without changing its chiral list.

\section{Every exceptional prime-field class is constructive}
The counts in Theorem~\ref{app:classification} are supplied by actual
matrices, rather than only a square criterion. Let $a,b,c$ now denote
chosen determinant-one traces, and put $d_0=a^2+b^2-4$.
For type $(5,3,4)$ choose $a=(-1\pm\sqrt5)/2$, $b=1$, $c=\sqrt2$;
for $(5,3,5)$ choose $a=c=(-1\pm\sqrt5)/2$, $b=1$;
for $(3,5,3)$ choose $a=c=1$, $b=(-1\pm\sqrt5)/2$.
Choose $U^2+V^2=d_0$, and for either square root $z$ of the associated
nonzero discriminant set
\[
 u=(U+a)/2,\quad w=(V-b)/2,\quad
 A=\begin{pmatrix}u&w+b\\w&a-u\end{pmatrix},\quad B=A^{-1}J,
 \qquad J=\begin{pmatrix}0&1\\-1&0\end{pmatrix}.
\]
With $R=(-bcV+zU)/d_0$, $S=(bcU+zV)/d_0$, put
\[
 C=\frac12\begin{pmatrix}R+c&S\\S&c-R\end{pmatrix}.
\]
The common trace theorem proves all determinant, trace and involution
identities, pair uniqueness and the exact number of roots. Its finite
intersection theorem proves that every such nonzero-root extension is
polytopal. Prime-field subgroup classification proves generation and
the nonzero invariant $2\operatorname{tr}(AC)-ac=z$ forbids reversal.
For $q\equiv3\pmod4$ these types require $q=p$ prime: the traces and
their quadratic extension all descend through an odd-degree ambient
field to $F_p$. This recovers the full scopes of source Theorems~6,
7, 9 and~11 without invoking their borrowed all-field quotient lists.

As a literal witness to the introductory omission, over $F_{571}$ take
\[
 A=\begin{pmatrix}0&481\\184&1\end{pmatrix},\quad
 B=\begin{pmatrix}481&1\\0&387\end{pmatrix},\quad
 C=\begin{pmatrix}113&411\\411&459\end{pmatrix}.
\]
These have type $(3,5,3)$; the two ends have order $60$, their
intersection order five, and the full group has order $93\,084\,420$.
The reversal obstruction is $\operatorname{tr}(AC)-1/2=550\ne0$.
Here $\sqrt5=24$, $571\equiv10\pmod{11}$ and the other introductory
clauses fail: $2$, $349$ and $229$ are nonsquares, the latter two being
$D_{55}$. The uniform proof above implies these conclusions, and the
separate complete permutation validation records the exact witness.
No new conjecture answer or historical-priority claim follows from this
version-specific correction.

\section{The earlier odd-composite-degree completion}
For odd coprime $e_1,e_2>1$, the common primitive-square theorem supplies
primitive $j_i\in F_{p^{e_i}}$ and a nonzero square root
$D^2=u^2v^2-4(u^2+v^2)$, $u=j_1+j_1^{-1}$, $v=j_2+j_2^{-1}$.
The common coprime-subfield theorem then gives the exact source matrices
\[
 \alpha=\frac12\begin{pmatrix}u&u+2\\u-2&u\end{pmatrix},\quad
 \beta=\frac12\begin{pmatrix}2+u&u\\-u&2-u\end{pmatrix},\quad
 \gamma=\frac1{2u}\begin{pmatrix}v(u-2)&D\\D&v(u+2)\end{pmatrix}.
\]
Their type is $((p^{e_1}-1)/2,p,(p^{e_2}-1)/2)$.
The scalar in $\gamma$ restores determinant one to the source's final
projective representative. The common proof verifies the shared
unipotent intersection before generation, then excludes every proper
maximal overgroup and every semilinear reversal. The norm-polynomial
proof includes equal and nested degrees, the complete integer reduction
and all eleven finite certificates. This material is fully reproduced
and paid in [U], with the prior phase267 acceptance recorded; it is not
an unpriced citation to an earlier proof bank and earns no new question credit.

\section{Exact source and evidence scope}
The source is the complete twenty-five-page arXiv:1606.08017v1 edition,
26 June 2016. No separate collation of its later journal edition is
claimed. We recover all its own principal classification and construction
clauses, with the exact corrected residue and characteristic qualifications.
The shared collision program also repairs its Lemma~5.12 calculation:
four of the twelve claimed trivial $A_4$-collision quotients are $A_5$,
but its intersection conclusion outside $L(4),L(5),L(19)$ is unchanged.
The source's general abelianization wording means a quotient, not a
subgroup; its finite cyclic application is unaffected. The signed-order
formula in Lemma~4.1 uses odd characteristic; in characteristic two
$-t=t$. The rank-three small identifications are the standard
$P(3)=S_4$, $P(2)=S_3$, correcting the printed sentence after Theorem~4.

Theorems~8 and~10 are quoted all-field Coxeter-quotient classifications
from Jones--Long and Jones--Long--Mednykh. They are not invoked in our
principal proof and their entire auxiliary all-field statements are not
newly rederived. Tables~1 and~2 are motivating finite census data,
including the $44$ objects at $q=169$; their full atlas computation is
not newly regenerated. The source map distinguishes these ancillary
items from every recovered principal clause and every explicit
construction used above. Its diagrams illustrate a replaced proof route.
No claim of literal all-assertions equivalence is made.

All three previous Topic E applications remain present, with their whole
common proof and complete required program. The additional required
finite premises and full polynomial JSON data are disclosed, with
validation and operational files separate. The nine new Lean algebra
statements and reused scoped certificates do not formalize classical
classification, character sums, finite GAP quotients or the polytope
correspondence. External proof, catalogue and graphic costs remain unknown.

\begin{thebibliography}{3}
\bibitem[U]{U} The human\_compression project,
\emph{Common theory for projective markings and rotations}, current complete common paper.
\bibitem{LMO} D. Leemans, J. Moerenhout and E. O'Reilly-Regueiro,
\emph{Projective linear groups as automorphism groups of chiral polytopes},
arXiv:1606.08017v1, 26 June 2016.
\bibitem{SW} E. Schulte and A. I. Weiss, \emph{Chiral polytopes},
Applied Geometry and Discrete Mathematics, DIMACS 4 (1991), 435--516.
\end{thebibliography}
\end{document}
